\section{Assessed Studies}

The identifiers below refer to the thirteen assessed studies. They do not
replace the record identifiers used in the search and selection files.
Page locators refer to the PDF version assessed, which can differ from
the pagination in the published proceedings.

\subsection{R01: Sage}
\noindent\textbf{Source and relevance.} \parencite{gan2021sage}. Published conference paper. LQ1, LQ3, LQ4. Sections 5-6, PDF pp. 10-11.

RPC graph and performance measurements support counterfactual diagnosis and resource actuation. Experiments inject resource interference and software faults in microservice applications.

\noindent\textbf{Interpretation limit.} The incomplete-instrumentation experiment retains per-tier latency and end-to-end measurements. Concurrent independent causes are harder than single causes.

\noindent\textbf{Appraisal.} Baselines require different instrumentation. Tested faults and platforms bound transfer. Repeated independent deployment variation is not established by the assessed sections.

\subsection{R02: TraceWeaver}
\noindent\textbf{Source and relevance.} \parencite{ashok2024traceweaver}. Published conference paper. LQ1, LQ3. Sections 4-6, PDF pp. 7-11.

Matches spans through timing distributions, call-graph constraints and joint optimisation. Evaluates benchmark applications and transformed Alibaba trace data, including load and dynamic call behaviour.

\noindent\textbf{Interpretation limit.} Request/response observations and test traces are required. Alibaba data are separated by call graph and inter-arrival times compressed. Reconstruction is not direct evidence of a cause.

\noindent\textbf{Appraisal.} Comparator assumptions differ. Production evaluation transforms sampled data. Useful ablations and load variation are reported, but this does not validate arbitrary cross-platform log correlation.

\subsection{R03: CIRCA}
\noindent\textbf{Source and relevance.} \parencite{li2022circa}. Published conference paper. LQ1, LQ3. Sections 5.1-5.4, Tables 3-4, PDF pp. 5-8.

Structural graph plus regression-based hypothesis testing ranks root-cause metrics. Oracle evaluation: 99 incidents, AC@1 0.404 and AC@5 0.763. Simulations assess graph and fault-strength variation.

\noindent\textbf{Interpretation limit.} Metric ranking differs from component diagnosis. Architecture knowledge and metric mapping are needed. Some baseline information and sampling assumptions differ.

\noindent\textbf{Appraisal.} Original baseline implementations do not all receive their original operator information or sampling rate. Real evaluation is one database context, with additional simulated robustness checks.

\subsection{R04: Causal RCA empirical comparison}
\noindent\textbf{Source and relevance.} \parencite{pham2024howfar}. Published conference paper. LQ1, LQ3. Sections 3-5, PDF pp. 3-9.

Compares causal discovery and RCA on synthetic data and microservice testbeds. Studies graph accuracy, diagnostic performance, runtime and input length, with repeated experiments.

\noindent\textbf{Interpretation limit.} CIRCA uses a PC-derived graph because its structural graph is unavailable. Limited fault types and benchmark systems constrain generalisation.

\noindent\textbf{Appraisal.} Extensive comparison and repetition, but adapted inputs and implementation choices prevent treating every result as a direct replication of the original method.

\subsection{R05: SIERRA}
\noindent\textbf{Source and relevance.} \parencite{lee2022sierra}. Published conference paper. LQ2, LQ3. Sections 4-7, Tables 2 and 5, PDF pp. 3-5 and 8-9.

Ranks anomalous SIEM windows using severity, direction and action context, and 62 count/rate/distribution features. Evaluates about 29.4 million firewall/IDS events with rolling weekly tests.

\noindent\textbf{Interpretation limit.} Ground truth is not exhaustive. Ten-minute windows and eight-week training periods may not suit a low-traffic PoC. Context features alone do not consistently improve detection.

\noindent\textbf{Appraisal.} Temporal evaluation and component comparisons are described. Partially labelled private enterprise data limit completeness and independent replication.

\subsection{R06: CICIDS2017 troubleshooting}
\noindent\textbf{Source and relevance.} \parencite{engelen2021troubleshooting}. Published conference paper. LQ2, LQ3. Sections III-IV, Tables I-II, PDF pp. 3-5.

Reconstructs dataset-generation steps and evaluates original and corrected flow data. Finds flow-construction and labelling defects and shortcut learning. Per-class outcomes change despite near-identical aggregate accuracy.

\noindent\textbf{Interpretation limit.} Findings concern one dataset and particular reconstruction and classifier choices, not every intrusion dataset.

\noindent\textbf{Appraisal.} Reconstruction and comparative experiments are inspectable. Dataset and model scope bound transfer beyond the studied benchmark.

\subsection{R07: Holmes}
\noindent\textbf{Source and relevance.} \parencite{chen2026holmes}. Preprint. LQ2, LQ3. Sections 3-4 and Table 2, PDF pp. 4-7 (arXiv v1).

Uses counters and sampled traffic for triage, then packet-derived evidence summaries and constrained reports. Representative replay windows include wrong attributions with high self-reported confidence.

\noindent\textbf{Interpretation limit.} One illustrative replay table is not an overall accuracy estimate. No comparative evidence establishes production effectiveness or safe automated repair.

\noindent\textbf{Appraisal.} Method is described, but the small illustrative evaluation lacks comprehensive comparison and repeated-run coverage. Claims about scalability are not established by this experiment.

\subsection{R08: RCACopilot}
\noindent\textbf{Source and relevance.} \parencite{chen2024rcacopilot}. Published conference paper. LQ3. Sections 5.1-5.6, Tables 2-4, PDF pp. 10-12.

Handlers collect diagnostics and retrieve historical incidents for root-cause categorisation. Uses 653 incidents with a 75/25 split. GPT-4 micro-F1 is 0.766, macro-F1 0.533. More context is not always better.

\noindent\textbf{Interpretation limit.} Private single-service labelled data and imbalanced categories limit transfer. Collection-module deployment scale must not be confused with equivalent deployment of LLM prediction.

\noindent\textbf{Appraisal.} Ablations and three rounds are reported. One service, private labels, model choice and historical similarity remain limitations.

\subsection{R09: RCAgent}
\noindent\textbf{Source and relevance.} \parencite{wang2024rcagent}. Published conference paper. LQ3. Sections 4-5, Tables 1-4, PDF pp. 5-6 (arXiv v2).

Controller and expert agents use bounded tools and observation snapshots. Evaluation uses 161 selected anomalous Flink jobs. Semantic and LLM-based scores, action validity and trajectory completion are distinct outcomes.

\noindent\textbf{Interpretation limit.} Tool and observation changes accompany the multi-agent design, so agent count is not isolated. LLM scoring and class-balanced selection affect interpretation. Proposed fixes are not executed repairs.

\noindent\textbf{Appraisal.} Engineer-reviewed annotations and leakage controls are described. Several architectural changes and surrogate output scores limit causal attribution to agent organisation.

\subsection{R10: TimeRAG}
\noindent\textbf{Source and relevance.} \parencite{park2026timerag}. Preprint. LQ1, LQ3. Sections 4-6, Tables 1-3, PDF pp. 5-7 (arXiv v1).

Aligns metric time series with textual representations and retrieves incident examples. Reports multiple-choice diagnosis accuracy separately from GPT-4o-judged report win rates. Some short-sequence results favour baselines.

\noindent\textbf{Interpretation limit.} Fixed-seed offline evaluation, augmented series and an LLM report judge limit operational conclusions. Report preference is not diagnosis accuracy.

\noindent\textbf{Appraisal.} Dataset and metric definitions are supplied. Augmentation, fixed seed and lack of live evaluation constrain robustness and transfer claims.

\subsection{R11: AIOpsLab}
\noindent\textbf{Source and relevance.} \parencite{chen2025aiopslab}. Published conference paper. LQ3, LQ4. Sections 3.2-3.6, Tables 2-4, PDF pp. 7-9.

Provides live fault-injection tasks for detection, localisation, RCA and mitigation. Compares agent behaviour, correctness, time, steps and tokens. More steps do not consistently improve performance.

\noindent\textbf{Interpretation limit.} The evaluated Flash implementation is simplified. Task definitions and selected problem sets affect comparisons. The reported selected-problem count and task-table denominators require care and are not pooled here.

\noindent\textbf{Appraisal.} Live evaluation and task definitions are useful. Reconstructed baselines, problem selection and reporting granularity qualify comparative conclusions.

\subsection{R12: STRATUS}
\noindent\textbf{Source and relevance.} \parencite{chen2025stratus}. Published conference paper. LQ3, LQ4. Sections 3-5, Tables 1-4, PDF pp. 5-9 (NeurIPS version).

Separates diagnostic and writer roles and uses reversible action stacks. GPT-4o solves 9/13 AIOpsLab mitigation tasks, versus 2/13 without retry and 3/13 with retry without undo.

\noindent\textbf{Interpretation limit.} Eight ITBench successes exploit non-persistent injected faults. Faithful undo and severity measurement are assumptions. Human approval support is described, but its effectiveness is not established by a user study.

\noindent\textbf{Appraisal.} Ablation supports undo-and-retry in the tested tasks. Small task sets, benchmark artefacts and differences in baseline provenance constrain generalisation.

\subsection{R13: SREGym}
\noindent\textbf{Source and relevance.} \parencite{clark2026sregym}. Preprint. LQ3, LQ4. Sections 2.5-3, Tables 2-5, PDF pp. 5-9 (arXiv v2).

Evaluates 90 live scenarios with three runs per agent-model configuration. Distinguishes diagnosis, mitigation and joint success. Added noise reduces diagnosis success for each configuration.

\noindent\textbf{Interpretation limit.} Different agents also have different models and tools. LLM-judge agreement is checked on a sample, not proof of an error-free oracle. One benchmark family cannot establish general production reliability.

\noindent\textbf{Appraisal.} Repeated runs and a human-labelled oracle-validation sample support appraisal. Architecture, tool and model differences confound a universal agent ranking.

\section{Quality-Appraisal Overview}

Table~\ref{tab:slr:quality} summarises the assistant judgements. S means
supported, P partly supported and N not supported. None of these
judgements establishes that the study can be reproduced independently.
The reasons and source locators above explain the material qualifications.

\begin{table}[htbp]
\centering\small
\caption{Separate quality judgements for the assessed evidence set.}
\label{tab:slr:quality}
\begin{tabular}{@{}llcccccc@{}}
\toprule
Study & Short name & Q1 & Q2 & Q3 & Q4 & Q5 & Q6 \\
\midrule
R01 & Sage & S & S & S & P & P & S \\
R02 & TraceWeaver & S & S & S & P & P & S \\
R03 & CIRCA & S & S & S & P & P & S \\
R04 & Causal RCA comparison & S & S & S & P & S & S \\
R05 & SIERRA & S & S & S & P & P & S \\
R06 & CICIDS2017 audit & S & S & S & S & P & S \\
R07 & Holmes & S & P & P & N & P & S \\
R08 & RCACopilot & S & S & S & P & P & S \\
R09 & RCAgent & S & S & S & P & P & S \\
R10 & TimeRAG & S & S & S & P & P & S \\
R11 & AIOpsLab & S & S & P & P & P & S \\
R12 & STRATUS & S & S & S & P & P & S \\
R13 & SREGym & S & S & S & P & S & S \\
\bottomrule
\end{tabular}
\end{table}
